% Оборот титульного листа: аннотация, адресат, как устроены примеры.
\pagestyle{empty}
\sffamily
\setlength{\parindent}{0pt}
\vspace*{\fill}

{\large\bfseries\color{GeoskanDarkBlue} Корягин С.\,А.}\par
\vspace{4pt}
{\color{GeoskanDarkBlue} Программирование дронов и создание автономных беспилотных миссий: учебное пособие по курсу «Программные средства БПЛА». --- Кудрово, \the\year. --- \pageref{LastPage}~с.}\par
\vspace{1cm}

{\bfseries\color{GeoskanDarkBlue} Аннотация}\par
\vspace{4pt}
Пособие вводит в программирование образовательных квадрокоптеров «Геоскан Пионер»: Пионера Базового (Lua и Python, библиотека \texttt{pioneer\_sdk}) и Пионера Мини~2 (Python на борту, библиотеки \texttt{pioneer\_sdk2}, \texttt{pioneer\_rknn}, ROS~2 и \texttt{pioneer\_swarm}, блоки Pioneer Bricks). Первая часть --- язык Lua, светодиодная индикация, таймеры и событийная модель автопилота, автономные миссии по точкам и параметрическим траекториям. Вторая часть --- Python: подключение, полёт, камера, визуальные маркеры ArUco, трекинг и групповой полёт Пионера Базового. Третья часть --- профессиональная навигация: калибровка камеры, карты маркеров, фильтрация измерений. Четвёртая часть --- Пионер Мини~2: программы на борту, нейросети на его нейропроцессоре и блочное программирование Bricks. Пятая часть --- ROS~2 и рой на Пионере: пакет \texttt{gs\_core} (библиотека \texttt{rclpy}: сервисы полёта, топики, параметры, узел-миссия) и Swarm Kit (библиотека \texttt{pioneer\_swarm}: соседи, сообщения, телеметрия, строй дронов). Программы запускаются в веб-симуляторе \texttt{simulator.sakoryagin.ru}: в нём есть обе модели дрона и режим \ui{Обучение} с пятью треками миссий и автоматической проверкой полёта. API симулятора совпадает с API настоящих дронов, а отличия названы в тексте.\par
\vspace{0.8cm}

{\bfseries\color{GeoskanDarkBlue} Для кого}\par
\vspace{4pt}
Для школьников и студентов, начинающих программировать роботов, и педагогов дополнительного образования. Предварительной подготовки по программированию не требуется; математика --- в объёме школьной тригонометрии.\par
\vspace{0.8cm}

{\bfseries\color{GeoskanDarkBlue} Примеры и проверка}\par
\vspace{4pt}
Исходный код всех примеров --- в папке \texttt{examples/} рядом с книгой; каждый файл подписан в тексте. Примеры для Lua и \texttt{pioneer\_sdk} автоматически прогоняются в симуляторе: взлетают, проходят точки и садятся. Примеры для Пионера Мини~2 (\texttt{examples/python-onboard/}) рассчитаны на борт дрона: их автоматически проверяют на имитации \texttt{pioneer\_sdk2} и \texttt{pioneer\_rknn} (по следу команд), примеры ROS~2 --- на эмуляции \texttt{rclpy} и узла \texttt{ros\_pioneer\_node}, а проекты Bricks сверяют с Python, который сгенерировала сама среда. Примеры роя проверяют на симуляторе с несколькими дронами. Полётные программы Мини~2, ROS~2 и роя на настоящем дроне авторами не запускались. Текст примеров можно копировать прямо из PDF.\par
\vspace{0.8cm}

{\footnotesize\color{GeoskanMuted} Версия \theversion\ · сборка \buildtimestamp\par}
\vspace*{\fill}
\clearpage
\pagestyle{scrheadings}
